

\newpage

%\footline={\hss\folio\hss}
%\pageno=1

\dspace
\subhead
3. Cohomologie et groupe fondamental 
\endsubhead

\subhead
3.1
\endsubhead
Soit $X$ un espace topologique s\'epar\'e et
connexe.
On note ${}_bP_a$ l'ensemble des classes
d'homotopie de chemins de $a$ \`a $b$, et
$\delta\gamma$ la composition des chemins
$_cP_b\times{} _bP_a\to {}_cP_a$.
Noter l'ordre des facteurs, choisi pour que si
$V$ est un syst\`eme local sur $X$ et que
$\gamma v\in V_b$ denote le transport de $v\in
V_a$ le long de $\gamma\in {}_bP_a$, on ait
$(\delta\gamma)v=\delta(\gamma v)$.

Nous ferons sur $X$ les hypoth\`eses (a) (b)
(c) suivantes.
(a) La th\'eorie classique du $\pi_1$
s'applique.
Par exemple: $X$ localement connexe par arcs
et localement simplement connexe.
(b) Le groupe fondamental est de type fini.
(c) Pour une projection $p\colon\, X^n\times
X^m\to X^n$ et un faisceau constant, $Rp_*$
commute au passage aux fibres.
Ces conditions sont remplies pour $X$
hom\'eomorphe au compl\'ement dans un poly\`edre
fini d'un sous-poly\`edre, par exemple pour
$X$ une vari\'et\'e alg\'ebrique complexe.

Fixons un corps $k$, le corps des
coefficients.
Soient $a\in X$, $\Lambda$ l'alg\`ebre de
groupe $k[\pi_1(X,a)]$ et $I$ son id\'eal
d'augmentation.
La donn\'ee d'un $\Lambda$-module \`a gauche
$V$, i.e. d'une repr\'esentation lin\'eaire de
$\pi_1(X,a)$ sur un $k$-espace vectoriel,
\'equivaut \`a celle d'un syst\`eme local 
$V\uptil$ sur $X$ de fibre $V$ en $a$.
Si on prend pour $V$ le $\Lambda$-module
$\Lambda$, on obtient le syst\`eme local de
fibre en $x$ l'espace vectoriel $k[{}_xP_a]$
de base ${}_xP_a$.
La structure de $\Lambda$-module \`a droite de
$\Lambda$ en fournit une sur ce syst\`eme
local $\Lamuptil$, i.e. une action \`a
droite de $\pi_1(X,a)$.
L'\'el\'ement $\gamma$ de $\pi_1(X,a)$ agit
sur ${}_xP_a$ par composition \`a droite.
Le syst\`eme local $(\Lambda/I^n)\uptil$ est
le quotient $\Lamuptil \otimes_\Lambda
\Lambda/I^n$ de $\Lamuptil$.

Le syst\`eme local $(\Lambda/I^n)\uptil$ est
extension it\'er\'ee des $n$ syst\`emes locaux
triviaux $I^p/I^{p+1}$ ($0\Le p<n$).
Sa fibre $(\Lambda/I^n)_a\!\!\!\uptil$ en $a$
est $\Lambda/I^n$ et l'unit\'e $1$ d\'efinit
$$
k\to (\Lambda/I^n)_a\!\!\!\uptil\colon\,
\lambda\longmapsto \lambda.1.
\tag3.1.1
$$
Le syst\`eme local $(\Lambda/I^n)\uptil$
est universel parmi les syst\`emes locaux $E$
extension it\'er\'ee de $n$ syst\`emes locaux
triviaux et munis de $k\to E_a$, car
$(\Lambda/I^n,1)$ est universel parmi les
$\Lambda$-modules $M$ tels que $I^nM=0$ munis
de $m_0\in M$.

\subhead
3.2
\endsubhead
Dans ce qui suit, nous laisserons de nombreux
signes ambigus.
Plusieurs conventions sont possibles, et
aucune ne nous semble clairement meilleure.
Elles n'importent gu\`ere, car elles
conduisent \`a des complexes isomorphe, avec
pour l'isomorphisme au pis l'ambigu{\umi}t\'e
d'un signe global.
Typiquement, changer de convention remplace un
complexe par un autre ayant les m\^{e}mes
composantes, et l'isomorphisme est donn\'e par
des signes d\'ependant de multidegr\'es.

\subhead
3.3
\endsubhead
Fixons $a,b\in X$ et consid\'erons dans $X^n$
les sous-espaces $b=t_1$,
$t_1=t_2,\dotsc,t_n=a$, not\'es $Y_0,\dots,Y_n$.
Pour $J\Subset\Delta_n:=\{0,1,\dotsc,n\}$,
soient $Y_J$ l'intersection des $Y_j$ ($j\in
J$) et $k_J$ le faisceau constant $k$ sur
$Y_J$, prolong\'e par z\'ero sur $X^n$.
Pour $J\Subset K$, on a $Y_K\Subset Y_J$ et on
dispose d'un morphisme de restriction $k_J\to
k_K$.
Si, pour $K=\{k_0,\dotsc,k_p\}$, les $k_\ell$
\'etant pris dans l'ordre croissant, et
$J=\{k_0,\dotsc,\khat_i,\dotsc,k_p\}$, on
l'affecte du signe $(-1)^i$, on obtient les
composantes de la diff\'erentielle d'un
complexe de faisceaux sur $X^n$:
$$
\underbark\to\oplus k_{\{j\}}\to\cdots\to
\oplusop_{\vert J\vert=p}
k_J\to\cdots\to k_{\Delta_n}.\tag3.3.1
$$
Noter que le dernier terme $k_{\Delta_n}$ est
nul sauf pour $n=0$ ou $a=b$, auxquels cas il
est r\'eduit \`a $k$ en $t_1=\ldots=t_n=a$.
Si on supprime le premier terme et qu'on
convient que les degr\'es vont de $0$ \`a $n$,
on obtient la r\'esolution \v{C}echiste
altern\'ee du faisceau constant $k$ sur la
r\'eunion $\cup Y_i$ des $Y_i$ pour le
recouvrement ferm\'e fini par les $Y_i$,
prolong\'ee par $0$ sur $X^n$.

D\'efinissons ${}_b\scrK_a\left<n\right>$,
ou simplement ${}_b\scrK_a$, comme \'etant le
complexe (3.3.1),
avec $k_{\Delta_n}$ omis, et les degr\'es
allant de $0$ \`a $n$.
Si $n>0$ et que $a\not=b$, c'est une
r\'esolution du prolongement par $0$ du
faisceau constant $k$ sur $X^n-\cup Y_i$, et
$$
\dbH^*(X^n,{}_b\scrK_a)=H^*(X^n\mod\,\cup
Y_i,k)\quad\text{(pour $n\not=0$,
$a\not=b$)}\tag3.3.2
$$

Si $a=b$, la derni\`ere diff\'erentielle
(3.3.1) fournit un morphisme de complexes
$$
\align
&{}_a\scrK_a\to k_{t_1=\cdots=
t_n=a}[-n],\tag3.3.3\\
\intertext{qui induit}
&\dbH^n(X^n,{}_a\scrK_a)\to k.\tag3.3.4
\endalign
$$

Si on r\'ep\`ete ces constructions en prenant
$a$ et $b$ comme param\`etres, on obtient un
complexe ${}_*\scrK_a$ sur $X\times X^n\times
X$, et, pour $p$ la projection $X\times
X^n\times X\to X\times X$, la restriction de
${}_*\scrK_*$ \`a $X^n\simover p^{-1}(b,a)$
s'identifie \`a ${}_b\scrK_a$.
Traitant seulement $b$ comme param\`etre, on
obtient de m\^{e}me ${}_*\scrK_a$ sur $X\times
X^n$.
Les $\dbH^i(X^n,{}_b\scrK_a)$ sont les fibres
du syst\`eme local $R^ip_*({}_*\scrK_*)$ sur
$X\times X$.

Nous allons, d'apr\`es Beilinson, v\'erifier la

\proclaim{Proposition 3.4}
(i) Pour $i<n$, $\dbH^i(X^n,{}_b\scrK_a)=0$.
\ (ii) Pour $b$ variable, le syst\`eme local
sur $X$ des $\dbH^n(X^n,{}_b\scrK_a)$, muni en
$b=a$ de (3.3.4), est le dual du syst\`eme
local $(\Lambda/I^{n+1})\uptil$ de 3.1, muni
du transpos\'e de (3.1.1).
\endproclaim

\demo{Preuve}
Pour $n=0$, $X^n$ est r\'eduit \`a un point,
${}_b\scrK_a$ est le faisceau constant $k$ en
degr\'e $0$ et (i) (ii) sont clairs.
Prouvons 3.4 par r\'ecurrence sur $n\Ge 1$;
l'hypoth\`ese de r\'ecurrence est la
validit\'e de 3.4 pour $m$ si $m<n$.

Soit $q$ la premi\`ere projection: $X^n\to X$.
Nous utiliserons la suite spectrale de Leray
pour $q$.
Ceci revient \`a \'ecrire que
$R\Gamma(X^n,\quad)=R\Gamma(X,Rq_*\quad)$ et
\`a filtrer $Rq_*$ par la filtration canonique
$\tau$.
Plus pr\'ecis\'ement, nous utiliserons une
filtration ayant une description uniforme pour
$a=b$ ou $a\not=b$, et qui revient \`a la
filtration canonique si $a\not=b$.

Un d\'evissage montre que la formation de
$Rq_*{}_b\scrK_a$ commute au passage aux
fibres.
Si on \'ecrit $X^n$ comme $X\times X^{n-1}$,
${}_b\scrK_a$ sur $X^n$ s'identifie au
c\^{o}ne$[-1]$ de
$$
\align
{}_*\scrK_a\left<n-1\right>\to{}_b\scrK_a\left<
n-1\right> &\quad\text{sur}\quad \{b\}\times
  X^{n-1}\tag3.4.1\\
\oplus~k_{\b{$\ssize t$}=a} &\quad\text{sur}\quad
\{a\}\times X^{n-1}
\endalign
$$
(complexe simple associ\'e au complexe double
de premi\`ere (resp. deuxi\`eme) colonne la
source (resp. but) de (3.4.1)).
\enddemo

On reconna{\hati}t \`a gauche les $k_J$ pour
$0\notin J$ et $J\not=\{1,\dotsc,m\}$, \`a
droite les $k_J$ d'une part pour $0\in J$, de
l'autre pour $J=\{1,\dotsc,n\}$.

Si $a\not=b$, la restriction de ${}_b\scrK_a$
\`a $q^{-1}(t)=X^{n-1}$ est donc: pour
$t\not=a,b$: ${}_t\scrK_a\left<n-1\right>$;
pour $t=b$: le c\^{o}ne$[-1]$ de
l'application identique de
${}_b\scrK_a\left<n-1\right>$; pour $t=a$: le
c\^{o}ne$[-1]$ de (3.3.3).
D'apr\`es l'hypoth\`ese de r\'ecurrence,
$R^iq_*{}_b\scrK_a=0$ si $i<n-1$.
Pour $i=n-1$, c'est un sous-faisceau du
syst\`eme local $(\Lambda/I^n)\uptil\upcheck$.
Il co{\umi}ncide avec
$(\Lambda/I^n)\uptil\upcheck$ si $t\not=a,b$,
est nul en $t=b$ et la fibre en a est le noyau
du transpos\'e de (3.3.1).
Pour $i\Ge n$, $R^iq_*{}_b\scrK_a$ est un
syst\`eme local sur $X-\{b\}$ prolong\'e par
$0$ sur $X$.
La suite spectrale de Leray v\'erifie donc
$E_2^{pq}=0$ pour $q\Le n-2$ ou $p=0$, et
$$
\align
&\dbH^i(X^n,{}_b\scrK_a)=0\quad\text{pour}\quad
 i<n\tag3.4.2\\
&\dbH^n(X^n,{}_b\scrK_a)=\dbH^1(X,R^{n-1}q_*
  {}_b\scrK_a)\quad\text{(pour $a\not=b$).}
\endalign
$$
Ce $\dbH^1$ est le groupe des extensions du
faisceau constant $k$ par
$(\Lambda/I^n)\uptil\upcheck$, trivialis\'ees
en $t=b$, et trivialis\'ees apr\`es avoir
pouss\'e par $(\Lambda/I^n)\upcheck\to k$ en
$t=a$.
Dualement, c'est le groupe des extensions
$\Etil$ de $(\Lambda/I^n)\uptil$ par $k$,
trivialis\'ees en $t=b$, et trivialis\'ees
au-dessus de $k\Subset\Lambda/I^n$ en $t=a$.

Traduisons en termes de $\Lambda$-modules.
A $\Etil$ correspond un $\Lambda$-module $E$,
extension de  $\Lambda/I^n$ par le module
d'augmentation $k$.
Il est annul\'e par $I^{n+1}$.
La donn\'ee additionelle en a est celle d'un
rel\`evement $\tilde{1}$ du g\'en\'erateur $1$
de $\Lambda/I^n$ dans $E$.
Soit $\alpha\colon\, I^n/I^{n+1}\to k$ la
restriction \`a $I^n/I^{n+1}$ du morphisme de
$\Lambda/I^{n+1}$ dans $E\colon\,
\lambda\mapsto\lambda \tilde{1}$.
L'extension $E$ est l'extension $E_\alpha$ de
l'extension $\Lambda/I^{n+1}$ de $\Lambda/I^n$
par $I^n/I^{n+1}$ en poussant par $\alpha$:
$$
\minCDarrowwidth{13pt}
\CD
0 @>>> I^n/I^{n+1} @>>> \Lambda/I^{n+1}
  @>>> \Lambda/I^n @>>> 0\\
@. @VV{\alpha}V @VVV @VVV\\
0 @>>> \dbQ @>>> E @>>> \Lambda/I^{n+1} @>>>
0.
\endCD
\tag3.4.3
$$

Pour $\alpha\colon\, I^n/I^{n+1}\to k$.
L'extension $E_\alpha\!\!\!\uptil$ de
$(\Lambda/I^n)\uptil$ par le syst\`eme local
trivial $k$ se d\'eduit de
$(\Lambda/I^{n+1})\uptil$ en poussant par
$$
\alpha\colon\text{ (syst\`eme local trivial
$I^n/I^{n+1})$}\to k.
$$
La trivialisation $\beta$ en $b$ fournit un
diagramme
$$
\minCDarrowwidth{13pt}
\CD
0 @>>> I^n/I^{n+1} @>>>
(\Lambda/I^{n+1})_b\!\!\!\uptil @>>>
(\Lambda/I^n)_b\!\!\!\uptil @>>> 0\\
@. @VV{\alpha}V @VVV @|\\
0 @>>> \dbQ @>>> (E_\alpha)_b\!\!\!\uptil
  &{\overset{\sssize\beta}\to\leftrightarrows} 
&(\Lambda/I^n)_b\!\!\!\uptil @>>> 0.
\endCD
\tag3.4.4
$$
De $\alpha$ et $\beta$ on d\'eduit une forme
lin\'eaire $\gamma$ sur
$(\Lambda/I^{n+1})_b\!\!\!\uptil$: \`a
$x\in(\Lambda/I^{n+1})_b\!\!\!\uptil$, d'image
$\xbar$ dans $(\Lambda/I^n)_b\!\!\!\uptil$,
attacher
$$
\gamma(x):=\image(x)-\beta(\xbar)\in\dbQ
\Subset(E_\alpha\!\!\!\uptil)_b.
$$
On a $\alpha=\gamma\vert(I^n/I^{n+1})$, et
$\gamma$ d\'etermine $\beta$.
Ceci fournit l'isomorphisme
$$
\dbH^n(X^n,{}_b\scrK_a)=
((\Lambda/I^{n+1})_b\!\!\!\uptil)\upvee
$$
annonc\'e.

Pr\'esent\'es un peu diff\'erement, ces
arguments continuent \`a s'appliquer pour
$a=b$.
L'image par $Rq_*$ du complexe ${}_b\scrK_a$,
c\^{o}ne$[-1]$ de (3.4.1) est encore un
c\^{o}ne$[-1]$.
On d\'efinit sa filtration croissante $\Fil$
en filtrant source et but par la filtration
canonique $\tau_{\Le}$.
Le gradu\'e $\Gr_i^{\Fil}$ s'envoie par un
quasi-isomorphisme sur le complexe
$$
R^iq_*{}_*\scrK_a\left<n-1\right>\to
\text{ (sa fibre en a)}_a\oplus
(k\text{ si }i=n-1)\tag3.4.5
$$
en degr\'es $0$ et $1$, translat\'e par
$[-i]$.
Pour $i<n-1$, (3.4.5) est nul.
Pour tout $i$, $\dbH^0((3.4.5))=0$.
La suite spectrale pour $R\Gamma(X,\quad)$ et
$\Fil$ donne donc
$$
\align
&\dbH^i(X^n,{}_b\scrK_a)=0\quad\text{pour
$i<n$}\tag3.4.6\\
&H^n(X^n,{}_b\scrK_a)=\dbH^1(X,(3.4.5)_{i=n-1}).
\endalign
$$
Ce $\dbH^1$ classifie les extensions du
faisceaux constant $k$ par
$(\Lambda/I^n)\uptil\upvee=R^{n-1}q_*{}_*\scrK_a
\left<n-1\right>$, trivialis\'ees apr\`es
avoir pouss\'e par le morphisme (3.4.5), et
les arguments donn\'es pour $a\not=b$ peuvent
\^{e}tre r\'ep\'et\'e.

Il reste, pour $a=b$, \`a identifier (3.3.4).
Le complexe (3.3.1) (degr\'es $0$ \`a $n+1$)
est le complexe simple associ\'e au complexe
triple de faisceaux sur $X^n=X\times X^{n-1}$:
$$
\CD
{}_*\scrK_a\left<n-1\right> @>>>{}_a\scrK_a
  \left<n-1\right>\text{ sur }q^{-1}(a)\\
@VVV @VVV\\
k[-(n-1)]_{{\b{$\ssize t$}}=a} @>{\sim}>> 
  k[-(n-1)]_{{\b{$\ssize t$}}=a}
\endCD
$$
et (3.3.4) est induit par le morphisme
$$
\CD
{}_*\scrK_a\left<n-1\right> @>>> 
  {}_a\scrK_a
  \left<n-1\right>\text{sur $\pr^{-1}(a)$}\\
@VVV \null\kern-3.0 cm @VVV\\
k[-(n-1)]_{{\b{$\ssize t$}}=a}
  @>>> \null\kern-3.0 cm 0
\endCD
\quad
\longrightarrow
\quad
\CD
0 @>>> \null\kern-25pt0\\
@VVV \null\kern-25pt@VVV\\
0 @>>> k[-(n-1)].
\endCD
$$

Filtrant comme pr\'ec\'edemment, on en
d\'eduit que (3.2.1) est le $\dbH^1$ du
morphisme de complexes
$$
\CD
(\Lambda/I^n)\upvee\uptil @>>>
(\Lambda/I^n)_a\!\!\!\uptil \oplus k_a\\
@VVV @VVV\\
0 @>>> k_a
\endCD
$$
(La fl\`eche verticale droite donnant le
morphisme est la diff\'erence entre
l'\'evaluation sur $k\Subset\Lambda/I^n$ et
l'identit\'e de $k_a$.)

Une extension de $k$ par
$(\Lambda/I^n)\upvee\uptil$, trivialis\'ee
apr\`es avoir poun\'e par (3.3.4), se dualise
en une extension de $(\Lambda/I^n)\uptil$ par
$k$ qui en $a$ est et trivialis\'ee, et
trivialis\'ee au-desus de $k\Subset\Lambda/I^n
=(\Lambda/I^n)_a\!\!\!\uptil$.
On lui attache la diff\'erence au-desus de $k$
entre ces deux trivialisations.
On laisse au lecteur le soin de v\'erifier que
cela revient \`a \'evaluer en $1$ une forme
lin\'eaire sur $\Lambda/I^{n+1}$.

\remark\nofrills{\bf 3.5 Remarque.}\enspace
Dans cette identification de 
$(k[{}_bP_a]\otimes_{\Lambda}\Lambda/I^{n+1})\upvee$
avec le groupe des extensions de
$(\Lambda/I^n)\uptil$ par $k$, munies de
trivialisations convenables en $a$ et $b$, le
sous-espace $(k[{}_bP_a]\otimes_\Lambda
\Lambda/I^n)\upvee$ est celui des extensions
triviales, avec la trivialisation \'evidente en
$a$ et une trivialisation en $b$.
L'injection de $(k[{}_bP_a]/I^n)\upvee\to
(k{}_aP_b]/I^{n+1})\upvee$ est donc le
$\dbH^n$ de
$$
\left({}_b\scrK_a\left<n-1\right>\text{ sur }
\{b\}\times X^{n+1}\right)[-1]\to
{}_b\scrK_a\text{ sur }X^n\tag3.5.1
$$
d\'eduit de la description du second membre
comme c\^{o}ne$[-1]$ de (3.4.1).
\endremark

\subhead
3.6
\endsubhead
Nous regarderons un objet simplicial
tant\^{o}t comme un foncteur contravariant sur la
cat\'egorie de tous les ensembles totalement
ordonn\'es finis non vides, tant\^{o}t comme un
foncteur sur la sous-cat\'egorie des
$\Delta_n=\{0,\dotsc,n\}$.
Cela revient au m\^{e}me (\'equivalence de
cat\'egories).
Si la valeur sur $\Delta_n$ d'un objet
simplicial $S$ est not\'ee $S_n$, on notera
$S_I$ sa valeur sur $I$ totalement ordonn\'e.
Elle est fonctorielle en $I$ et si $I$ a $n+1$
\'el\'ement, l'unique isomorphisme de $I$ avec
$\Delta_n$ fournit un isomorphisme de $S_I$
avec $S_n$.

Le morphismes $\Delta_n\to\Delta_1$ sont les
$$
\varphi_i\colon\, 0,\dotsc,i-1\longmapsto
0;\,\,
i,\dotsc,n\longmapsto 1\qquad
(0\Le i\Le n+1).\tag3.6.1
$$
Nous identifierons par cette num\'erotation
$X^{\Hom(\Delta_n,\Delta_1)}$ \`a $X\times
X^n\times X$.
L'ensemble simplicial
$\Delta_n\mapsto\Hom(\Delta_n,\Delta_1)$
contient l'ensemble simplicial
$\Delta_n\mapsto\Hom(\Delta_n,\partial\Delta_1)=
\{\varphi_0,\varphi_{n+1}\}$.
L'espace cosimplicial $\Delta_n\mapsto
X^{\Hom(\Delta_n,\Delta_1)}$ se projette donc sur
l'espace cosimplicial $\Delta_n\mapsto
X^{\Hom(\Delta_n,\partial\Delta_1)}=
X\times X$.
Prenant la fibre en $(b,a)$, on obtient un
espace cosimplicial $\Delta_n\mapsto X^n$.
Les $\partial_i\colon\, X^{n-1}\to X^n$
($0\Le i\Le n$) sont les
$$
\alignat2
t_1,\dotsc,t_{n-1}\longmapsto
  &~b,t_1,\dotsc,t_{n-1} &\quad &(i=0)\\
&~t_1,\dotsc,t_i,t_i,\dotsc,t_{n-1} 
  &\quad &(i\not=0,n)\\
&~t_1,\dotsc,t_{n-1},a &\quad &(i=n)\,\,,
\endalignat
$$
d'images les $Y_i$ de 3.3.
Plus g\'en\'eralement, pour $J\Subset\{
0,\dotsc,n\}$ de compl\'ement $I$, suppos\'e
non vide, $p=\vert I\vert-1$ et $\varphi_I$
l'inclusion de $I$, le morphisme cosimplicial
$\varphi_I\colon\, X^p\to X^n$ identifie $X^p$
avec $Y_J\Subset X^n$.
Pour $J$, $J'$ de compl\'ements $I$, $I'$ et
$p=\vert I\vert-1$, $p'=\vert I'\vert-1$, si
$J\Subset J'$ et que $\varphi_{II'}$ est
l'inclusion de $I'$ dans $I$, le morphisme
cosimplicial $\varphi_{II'}\colon\, X^{p'}\to
X^p$ s'identifie \`a l'inclusion de $Y_{J'}$
dans $Y_J$.
Noter que pour $J=\{j_0,\dotsc,\Hatj_k,\dotsc,j_p
\}\Subset J'=\{j_0,\dotsc,j_p\}$ donnant lieu
\`a $I'=\{i_0,\dotsc,{\Hati}_\ell,\dotsc,i_q\}
\Subset I=\{i_0,\dotsc,i_q\}$ (les $i$ et $j$
sont pris dans l'ordre croissant), on a
$p+q=n$, $k+\ell=i_\ell=j_k$, et le signe
$(-1)^k$ que nous avons employ\'e en 3.3 peut
se r\'ecrire
$$
(-1)^k=(-1)^\ell\cdot
(-1)^{j_k}=\eps(J)(-1)^\ell\eps(J')
$$
pour
$$
\eps(J)=\prod\limits_{j\in J}(-1)^j.\tag3.6.2
$$

\subhead
3.7
\endsubhead
Choisissons des complexes $S_m^*$ qui
calculent la cohomologie des $X^m$, et soient
fonctoriels pour les morphismes cosimpliciaux
entre les $X^m$.
Plus pr\'ecis\'ement: on prend pour $S_m^*$ le
complexe des sections globales d'une
r\'esolution fonctorielle $\scrS_m^*$ du
faisceau constant $k$ sur $X^m$, telle que
$\Gamma(X^m,\scrS_m^*)\to
R\Gamma(X^m,\scrS_m^*)$ soit un
quasi-isomorphisme.
A $m$ variable, les $S_m^*$ forment un
syst\`eme simplicial de complexes.
Les $\scrS_m^*$ fournissent une r\'esolution
de ${}_b\scrK_a\left<n\right>$ sur $X^n$:
avec les notations de 3.6, r\'esoudre $k_J$
par $\varphi_{I*}\scrS_I^*$ et utiliser la
fonctorialit\'e de $\scrS_I^*$ en $I$.

Un {\it syst\`eme de coefficients} $c$ sur le
simplexe standard $\Delta_n$, \`a valeurs dans
une cat\'egorie additive, attache \`a chaque
facette $F\Subset\Delta_m$ un objet $c(F)$,
contravariant en $F$ pour les morphismes
d'inclusion.
Il d\'efinit un complexe de cha{\hati}nes
$$
C_p(\Delta_n,c)=\bigoplus\limits_{\vert F\vert
=p+1}c(F).
$$
Dualement, un cosyst\`eme: $c(F)$ covariant en
$F$, d\'efinit un complexe de cocha{\hati}nes.

Si $s_{\smbul}$ est un objet simplicial, on
note encore $C_*(\Delta_n,s_{\smbul})$ le
complexe $C_*(\Delta_n,c)$ pour $c(F):=s_F$.
Dualement pour $s^{\smbul}$ cosimplicial.

Le complexe
$$
\Gamma(X^n,\text{ r\'esolution de
${}_b\scrK_a$ d\'efinie par $\scrS_{\smbul}$}),
$$
calcule l'hypercohomologie de ${}_b\scrK_a$.
Il s'identifie (avec les signes (3.6.2)) au
complexe simple associ\'e au complexe double
($=$ complexe de complexes)
$$
C_*(\Delta_n,S_{\smbul}^*),
$$
translat\'e par $[-n]$, et 3.4 (ii) se
r\'ecrit
$$
(k[{}_b P_a]\otimes_\Lambda
\Lambda/I^{n+1})\upvee=H^0C_*(\Delta_n, S_{\smbul}^*).
\tag3.7.1
$$

\example\nofrills{3.8 Exemple:}\enspace
prenons pour $S_m^*$ le complexe des cocha{\hati}nes
singuli\`eres de $X^m$.
Il ne rentre pas tout \`a fait dans le cadre
pr\'ec\'edent, mais on peut l'y ramener: si
$\scrS_m^*$ est le complexe des cocha{\hati}nes
singuli\`eres localis\'ees sur $X^m$ (le
faisceau associ\'e au pr\'efaisceau des
cocha{\hati}nes singuli\`eres),
$S_m^*\to\Gamma(X^m,\scrS_m^*)$ est un
quasi-isomorphisme.
Le complexe $S_m^*$ des cocha{\hati}nes
singuli\`eres est le dual du complexe $S_*^m$
des cha{\hati}nes singuli\`eres.
A $m$ variable, $S_*^m$ est cosimplicial en $m$,
et $C_*(\Delta_n,S_{\smbul}^*)$ est le dual de
$C^*(\Delta_n,S_*^{\smbul})$.
On a donc
$$
k[{}_b P_a]\otimes_\Lambda
\Lambda/I^{n+1}=H^0C^*(\Delta_n,S_*^{\smbul})
\tag3.8.1
$$

Pour $X=[0,1]$, $a=0$ et $b=1$, ${}_b P_a$ est
r\'eduit \`a un \'el\'ement et chaque
$k[{}_bP_a]\otimes_\Lambda \Lambda/I^{n+1}$
est r\'eduit \`a $k$.
Au membre de gauche, les
$C^*(\Delta_n,S_*^{\smbul})$ forment un
syst\`eme projectif, de morphismes de
transition d\'eduit, apr\`es dualisation, de
(3.5.1\quad):
$$
C^*(\Delta_n,S_*^{\smbul})\to
C^*(\Delta_{n-1},
S_*^{\smbul})
$$
est induit par $\partial_0\colon\,\Delta_{n-1}
\hookrightarrow\Delta_n$.
\endexample

Pour qu'un cocycle de degr\'e zero de
$C^*(\Delta_n,S_*^{\smbul})$ corresponde au
chemin de $0$ \`a $1$, il faut et il suffit
que sa valeur en $\{n\}=\partial_0^n(\Delta_0)\Subset
\Delta_n$ soit $1\in S^0$, le complexe des
cha{\hati}nes singuli\`eres du point.
On a $[0,1]=\vert\Delta_1\vert$.
Soit $\sigma_m$ le simplexe
$\vert\Delta_m\vert
\to\vert\Delta_1\vert^m=[0,1]^m$ de
coordonn\'ees les d\'eg\'en\'erescences
$\varphi_i$ ($1\Le i\le m$) de (3.6.1).
Son image est $\{\underbart\vert 1\Ge
t_1\Ge\ldots\Ge t_m\Ge 0\}$.
Comme cocycle, on peut prendre celui qui \`a
chaque face de dimension $p$ de $\Delta_n$
attache $\pm\sigma_p$, le signe \'etant $+$
pour $p=0$ et d\'ependant des conventiones de
signes d\'efinissant
$C^*(\Delta_n,S_*^{\smbul})$ pour $p>0$.

Par fonctorialit\'e en $X$, on en d\'eduit
que, pour $X$, l'image dans
$H^0C^*(\Delta_n,S_*)$ d'un chemin
$\gamma\colon\, [0,1]\to X$ de $a$ \`a $b$ est
repr\'esent\'ee par le cocycle qui \`a chaque
face de dimension $p$ de $\Delta_n$ attache
$\pm\gamma\circ\sigma_p$:
$\vert\Delta_p\vert\to[0,1]^p\to X^p$, le
signe \'etant comme ci-desssus.

\subhead
3.9
\endsubhead
La litt\'erature donne une autre recette pour
d\'eduire $(k[{}_bP_a]\otimes_\Lambda
\Lambda/I^{n+1})\upvee$ de l'objet simplicial
$S_{\smbul}^*$ de la cat\'egorie des complexes
consid\'er\'e en 3.7:

\medskip\noindent
{\sans (a)}\enspace
prendre le complexe normalis\'e
$NS_{\smbul}^*$ de cet objet simplicial;

\smallskip\noindent
{\sans (b)}\enspace
prendre le tronqu\'e $\sigma_{\Ge -n}
NS_{\smbul}^*$ qui ne garde que les
$NS_m^*$ pour $m\Le n$;

\smallskip\noindent
{\sans (c)}\enspace
prendre le complexe simple associ\'e au complexe
double $\sigma_{\Ge -n}NS_{\smbul}^*$ et son
$H^0$.

\medskip
L'\'equivalence entre les deux constructions
r\'esulte du lemme suivant.
Rappelons qu'une {\it cat\'egorie karoubienne}
est une cat\'egorie additive dans laquelle
tout endomorphisme idempotent admet un noyau,
et que l'\'equivalence de Dold-Puppe:
(objet simplicial
$S_{\smbul}$)$\to$ (complexe $NS_{\smbul}$ \`a
degr\'es $\Le 0$) est disponible dans une
telle cat\'egorie ($NS_k$ est facteur
direct de $S_k$)

\proclaim{Proposition 3.10}
Si $S_{\smbul}$ est un objet simplicial d'une
cat\'egorie karoubienne, le complexe 
$C_*(\Delta_n,S_{\smbul})$ et le
tronqu\'e normalis\'e $\sigma_{\Ge
-n}NS_{\smbul}$ sont fonctoriellement homotopes.
\endproclaim

Nous v\'erifierons par la m\'ethode des
mod\`eles acycliques l'\'enonc\'e dual, pour
un objet cosimplicial $S^{\smbul}$,
$\sigma_{\Le n}NS^{\smbul}$ et
$C^*(\Delta_n,S^{\smbul})$.

\demo{Preuve}
Nous noterons $C_{\dbZ}$ la cat\'egorie
pr\'eadditive ayant les m\^{e}mes objets
qu'une cat\'egorie $\scrC$, pour laquelle
$\Hom(A,B)$ est le groupe ab\'elien librement
engendr\'e par $\Hom_{\scrC}(A,B)$.
Nous noterons $\scrC_{\kar}$ son enveloppe
karoubienne, obtenue par adjonction formelle
de sommes finies et de facteurs directs
noyaux d'endomorphismes idempotents.
Le foncteur contravariant
$$
X\longmapsto h_{\dbZ}^X\colon\,
\scrC_{\kar}\to\HOM(\scrC,Ab),\quad 
h_{\dbZ}^X(Y)=\Hom_{\scrC_{\kar}}(X,Y)
$$
est pleinement fid\`ele.
Il induit une \'equivalence de $\scrC_{\kar}$
avec la sous-cat\'egorie de la cat\'egorie
ab\'elienne $\HOM(\scrC,Ab)$ d'objets les
facteurs directs de sommes finies de
foncteurs
$$
Y\longmapsto\text{ groupe ab\'elien librement
engendr\'e par $\Hom_{\scrC}(X,Y)$},
$$
pour $X$ dans $\scrC$.
Ces foncteurs sont des objets projectifs de la
cat\'egorie ab\'elienne $\HOM(\scrC,Ab)$, car
pour $C$ dans $\scrC$
$$
\Hom(h_{\dbZ}^C,F)=F(C)
$$
(variante du lemme de Yoneda).
\enddemo

Si $\Delta$ est la cat\'egorie des $\Delta_n$,
l'objet cosimplicial $(\Delta_n)$ de
$\Delta_{kar}$ est universel en ce sens que,
pour $\scrA$ une cat\'egorie karoubienne,
$F\mapsto F((\Delta_n))$ est un \'equivalence
$$
\text{(foncteurs $\Delta_{\kar}\to\scrA)\to$
(objets cosimpliciaux de $\scrA$)}
$$
Pour v\'erifier l'\'enonc\'e $(3.10)^*$ dual de
3.10, et montrer que l'\'equivalence
d'homotopie peut \^{e}tre donn\'ee par des
formules universelles, il suffit de v\'erifier
$(3.10)^*$ dans le cas universel: pour l'objet
$(\Delta_n)$ de $\Delta_{\kar}$.
Si $\Delta_{\Le n}$ est la sous-cat\'egorie de
$\Delta$ d'objets les $\Delta_m$ pour $m\Le
n$, tant $\sigma_{\Ge-n}N\Delta_{\smbul}$ que
$C_*(\Delta_n,\Delta_{\smbul})$ sont dans
$(\Delta_{\Le n})_{\kar}$, et c'est dans
$(\Delta_{\Le n})_{\kar}$ que nous
travaillerons.

Appliquons le foncteur $X\mapsto h_{\dbZ}^X$
\`a $\sigma_{\Le n}N\Delta^{\smbul}$ et \`a
$C^*(\Delta_n,\Delta^{\smbul})$.
On obtient deux complexes d'objets projectifs
de la cat\'egorie ab\'elienne
$\HOM(\Delta_{\Le n},Ab)$.
Nous montrerons qu'il sont l'un et l'autre une
r\'esolution projective du foncteur constant
$\dbZ$, donc sont homotopes, et on conclut par la
pleine fid\'elit\'e de $X\mapsto h_{\dbZ}^X$.
Si on \'evalue en $\Delta_p$ ($p\Le n$), on
obtient les complexes $\sigma_{\Ge
-n}NS_{\smbul}$ et $C_*(\Delta_n,S_{\smbul})$
de 3.10 pour $S_{\smbul}$ le complexe
cha{\hati}nes simpliciales de $\Delta_p$.
Ce sont

\medskip\noindent
{\sans (a)}\enspace
le complexe des cha{\hati}nes non
d\'eg\'en\'er\'ees de $\Delta_p$, \'egal \`a
son tronqu\'e car $p\le n$.

\smallskip\noindent
{\sans (b)}\enspace
le complexe des cha{\hati}nes combinaisons
lin\'eaires de
$\Delta_q\to\Delta_n\times\Delta_p$, avec la
premi\`ere projection injective.

\medskip
Ces deux complexes sont augment\'es vers $\dbZ$,
fonctoriellement en $\Delta_p$.
Puisque $\vert\Delta_p\vert$ est contractile,
et que le complexe (a) calcule son homologie,
il est une r\'esolution de $\dbZ$.
Le complexe (b) est le complexe des
cha{\hati}nes non d\'eg\'en\'er\'ees du
sous-polytope simplicial suivant de
$\vert\Delta_n\times\Delta_p\vert$:

\parindent=40pt
\medskip\noindent
\Item{(3.10.1)}
la r\'eunion, pour $\varphi\colon\,
\Delta_n\to\Delta_p$, des graphes
$\Gamma_\varphi$ des morphismes
$\vert\varphi\vert\colon\,\vert\Delta_n\vert\to
\vert\Delta_p\vert$.

\parindent=20pt
\medskip
Pour v\'erifier que le complexe (b) est une
r\'esolution de $\dbZ$, il suffit donc de
v\'erifier le

\proclaim{Lemme 3.10.2}
L'espace (3.10.1) est contractile.
\endproclaim

\demo{Preuve}
Ordonnons lexicographiquement l'ensemble des
applications croissantes
$$
\varphi\colon\,\{0,\dotsc,n\}\to
\{0,\dotsc,p\}.
$$
Montrons que pour tout $\varphi\not=0$,
$$
\Gamma_\varphi\CAP\bigcup\limits_{\varphi' < \varphi}
\Gamma_{\varphi'}
$$
est une union non vide de faces du simplexe
$\Gamma_{\varphi}$, mais non de toutes les faces
de $\Gamma_\varphi$, de sorte que
$\bigcup\limits_{\varphi'<\varphi}\Gamma_{\varphi'}$
est r\'etracte par d\'eformation de
$\bigcup\limits_{\varphi'\Le
\varphi}\Gamma_{\varphi'}$.

Quels que soient $\varphi'$ et $\varphi$, si
$F\Subset\{0,\dotsc,n\}$ est l'ensemble des
$i$ tels que $\varphi(i)=\varphi'(i)$, et
$\vert F\vert$ la facette correspondante
$\Delta_m$, $\Gamma_\varphi\CAP
\Gamma_{\varphi'}$ est le graphe de la
restriction de $\vert\varphi\vert$ ou
$\vert\varphi'\vert$ \`a $\vert F\vert$.

Identifions $\Gamma_\varphi$ \`a
$\vert\Delta_m\vert$ par la projection sur
$\Delta_m$.
Si $\varphi'<\varphi$, soit $i$ minimal tel
que $\varphi'(i)<\varphi(i)$.
Si $\varphi''$ est d\'efini par
$\varphi''(j)=\varphi'(j)$ (resp.
$\varphi(j)$) pour $j\le i$ (resp. $j>i$), on
a $\varphi''<\varphi$ et
$\Gamma_\varphi\CAP\Gamma_{\varphi'}\Subset
\Gamma_\varphi\CAP\Gamma_{\varphi''}$, qui
s'identifie \`a la face d'indice $i$ de
$\vert\Delta_m\vert$.
L'intersection
$\Gamma_\varphi\CAP\bigcup\limits_{\varphi'<\varphi}
\Gamma_{\varphi'}$ est donc r\'eunion non vide
de faces.

Pour que la face d'indice $i$ soit obtenue, il
faut et il suffit que $\varphi(i)\not=0$, et
que, si $i\not=0$, on ait
$\varphi(i-1)<\varphi(i)$.
Parce que $n\Ge p$, toute les faces ne sont
pas obtenues: soit $\varphi(i)=\varphi(i+1)$
pour un $i$, soir $n=m$, $\varphi$ est
l'identit\'e, et $\varphi(0)=0$.
\enddemo

\subhead
3.11
\endsubhead
Si $X$ est une vari\'et\'e diff\'erentiable,
que $k=\dbR$ et que, dans 3.7, on calcule la
cohomologie de $X^n$ \`a l'aide du complexe de
de Rham, on obtient un objet simplicial
$\Omega^p(X^n)$ de la cat\'egorie des
complexes (degr\'e diff\'erentiel gradu\'e $p$,
degr\'e simplicial $n$).
Soit $N\Omega^*(X^*)$ son normalis\'e, et
$\bfs N\Omega^*X^*$ le complexe simple
associ\'e, dont la composante de degr\'e $k$
est $\oplusop_{p-n=k}N\Omega^p(X^k)$.
C'est un proche cousin du complexe des
int\'egrales it\'er\'ees de Chen.
Nous nous proposons d'expliquer la relation
entre les deux.

Une $q$-forme $\beta$ sur l'espace
${}_b\Omega_a(X)$ des chemins $C^\infty$ de
$a$ \`a $b$ est pour Chen la donn\'ee, pour
toute vari\'et\'e $U$ et toute famille
$C^\infty$ $^\colon\, U\to {}_b\Omega_a(X)$
de chemins de $a$ \`a $b$ param\'etr\'ee par
$U$, d'une $q$-forme (not\'ee
$\gamma^*\beta$) sur $U$, cette donn\'ee
v\'erifiant que pour $f\colon\, V\to U$ on
a $(\gamma f)^*\beta=f^*(\gamma^*\beta)$.
A une $p$-forme sur $X^n$, Chen attache comme
suit une $(p-n)$-forme sur ${}_b\Omega_a(X)$:

\medskip\noindent
{\sans (a)}\enspace
r\'ep\'etant avec param\`etres la construction
de 3.8, on attache \`a une famille $\gamma$ de
chemins param\'etr\'ee par $U$ un morphisme
$$
\gamma_n=\gamma\circ\sigma_n\colon\,
U\times\vert\Delta_n\vert\to
U\times[0,1]^n\to X^n;
$$

\noindent
{\sans (b)}\enspace
\`a la forme $\alpha$ sur $X^n$ on attache
l'int\'egrale de $\gamma_n^*\alpha$ le long
des fibres de $U\times\vert\Delta_n\vert\to U$.
Plus pr\'ecis\'ement, on attache
$$
\left<\alpha\right>=(-1)^{n(n+1)/2}\int_{U\times
\vert\Delta_n\vert/U}\gamma_n^*\alpha,
$$
o\`u l'int\'egrale le long des fibres est
d\'efinie de sorte que pour $\beta'$ sur
$U$ et $\beta''$ sur $\vert\Delta_n\vert$,
on ait 
$$
\int_{U\times\vert\Delta_n\vert/U}
\pr_1^*\beta'\w \pr_2^*\beta''=\beta'.
\int_{\vert\Delta_n\vert}\beta''.
$$

Le signe est choisi pour avoir, si $\alpha$
est de degr\'e $p$
$$
d\left<\alpha\right>=\left<d\alpha\right>+
(-1)^p\left<\sum(-1)^ij_i^*\alpha\right>\,\,:
$$
compatibilit\'e de la diff\'erentielle
ext\'erieure sur ${}_b\Omega_a(X)$ avec la
diff\'erentielle $D$ du complexe simple
associ\'e au complexe double des
$\Omega^*(X^n)$ (pour * le premier degr\'e).

Pour la diff\'erentielle transpos\'ee de
$\Sing_*X^*$, si $\gamma$ est un chemin de $a$
\`a $b$, c'est $\sum(-1)^p\gamma\circ\sigma_p$
qui est un cycle.

Chen remplace les $\Omega^*X^n$ par le
syst\`eme simplicial de sous-complexes
$\otimesop^{n}\Omega^*(X)$.
Le {\it complexe des int\'egrale it\'er\'ees}
de Chen est l'image par $\left<\quad\right>$
de la somme des $\otimesop^{n} \Omega^*(X)$.
Cette image est form\'ee des sommes de 
$$
\left<\alpha_1\vert\ldots\vert\alpha_n\right>:=
\left<\pr_1^*\alpha_1\w\ldots\w\pr_n^*\alpha_n
\right>.
$$

Les d\'eg\'en\'erescences $s_i\colon\, X^n\to
X^{n-1}$ sont les projections $(x_1\ldots
x_n)\to
(x_1\ldots\xhat_i\ldots x_n)$ et le complexe
normalis\'e de $\otimesop^{n}\Omega^*X$, dans
sa version ``diviser par les $\im(s_i)$'' est
$$
N\otimesop^{n}\Omega^*X=\otimesop^{n}(\Omega^*X/
\text{constantes}).\tag3.11.1
$$
Par K\"unneth, les $\otimesop^{n}\Omega^*X$
sont quasi-isomorphes aux $\Omega^*X^n$, et
donc les $N\otimesop^{n}\Omega^*X$ aux
$N\Omega^*X^n$.

Si un des $\alpha_i$ est de degr\'e $0$,
$\left<\alpha_1\vert\ldots\vert\alpha_n\right>$
est nul et $\left<\quad\right>$ induit donc
$$
\left<\quad\right>\colon\,
N\otimesop^{n}\Omega^*X\to\text{complexe des
int\'egrales it\'erees.}
$$
Les deux membres sont filtr\'es par $n$; par
les $\oplusop_{m\Le n}N\otimesop^{m}\Omega^*X$
et leurs images.
Chen (1973) (Lemme 4.3.1) prouve 
que, pour cette filtration $\Fil$ du complexe
des int\'egrales it\'eries, $\Gr_n^{\Fil}$ est
le quotient $\otimesop^{n}(\tau_{\Ge 1}\Omega^*X)$
de $\otimesop^{n}\Omega^*X$.
Le complexe $\tau_{\Ge 1}\Omega^*X$ est le
quotient de $\Omega^*X$ par $\Omega^0X$ et son
image par $d$.
La vari\'et\'e $X$ \'etant connexe, le
morphisme de $(\Omega^*X/\text{constantes})$
dans $\tau_{\Ge 1}\Omega^*X$ est un
quasi-isomorphisme, et les
$$
\Fil_nN\otimesop^{n}\Omega^*X
\to\Fil_n(\text{int\'egrales it\'er\'ees})
$$
sont donc des quasi-isomorphismes.

\remark{Remarque}
On peut donner du complexe des int\'egrales
it\'er\'ees une description directe, en terme
de l'alg\`ebre gradu\'ee $\Omega^*(X)$ et des
morphismes d'\'evaluation en $a$ et $b$, sans
invoquet
les formes diff\'erentielles sur
${}_b\Omega_a$: l'image de $\oplusop_{m\Le
n}\otimesop^{m}\Omega_X^*$ est son quotient par
les $\alpha_1\otimes\ldots\otimes\alpha_m$
tels qu'un $\alpha_i$ soit de degr\'e $0$, et
leurs diff\'erentielles (construction bar
r\'eduite de Chen 1976).
\endremark

\subhead
3.12
\endsubhead
Soit $X$ une vari\'et\'e alg\'ebrique lisse
connexe sur un corps $K$ et $a,b\in X(K)$.
Comme dans le cas topologique 3.6, le sch\'ema
cosimplicial $X^{\Hom(\Delta_n,\Delta_1)}$
s'envoie dans le sch\'ema cosimplicial constant
$X\times
X=X^{\Hom(\Delta_n,\partial\Delta_1)}$ et,
prenant la fibre en $(b,a)\in X\times X$, on
obtient un sch\'ema cosimplicial
$\Delta_n\mapsto X^n$.
Regardons-le comme un objet cosimplicial de la
cat\'egorie additive $\SmCor(K)$.
Quel que soit $n$, on peut lui attacher le
complexe
$$
C^*(\Delta_n,X^*)[n]
$$
(3.7) de $\SmCor(K)$, ou, dans l'enveloppe
karoubienne $\SmCor(K)_{\kar}$ de $\SmCor(K)$,
le complexe homotope $\sigma_{\Le n}NX^*$.
Ces complexes fournissent des objets,
naturellement isomorphes d'apr\`es 3.10, de la
cat\'egorie triangul\'ee motivique $\DM(K)$.
Notons les ${}_b\Omega_a^{[n]}(X)$.

Appliquons \`a ${}_b\Omega_a^{[n]}(X)$ le
foncteur ``r\'ealisation''  cohomologique,
suivi de $H^0$.
D'apr\`es 3.4, pour tout plongement complexe
$\sigma$ de $K$, dans la r\'ealisation de
Betti $H_\sigma$ correspondante (coefficients
rationnels), on obtient
$(\dbQ[{}_bP_a]\otimes_\Lambda
\Lambda/I^{n+1})\upvee$ pour $X(\dbC)$.

Supposons que $X$, comme objet de $\DM(K)$,
soit de Tate mixte, i.e. dans $\DMT(K)$.
Les ${}_b\Omega_a^{[n]}(X)$ sont alors
\'egalement dans $\DMT(K)$.
Si $K$ est un corps de nombres, ou plus
g\'en\'eralement quand on dispose de la
conjecture d'annulation de Beilinson-Soul\'e,
on peut, apr\`es tensorisation avec $\dbQ$,
appliquer $H^0$ et dualiser pour obtenir un
objet de la cat\'egorie $\MT(K)$ des motifs de
Tate mixtes sur $K$.

Si on travaille avec $\sigma_{\Le n}NX^*$, les
formules classiques en topologie (voir 
Wojtkowiak (1993)) fournissent, apr\`es
passage \`a la limite inductive en $n$, une
structure d'alg\`ebre commutative dans la
cat\'egorie des $\Ind$-objets de $\MT(K)$, et
$$
{}_bA_a(X):=\colim_n(H^0\sigma_{\Le
n}NX^*)\upvee
$$
m\'erite le nom d'alg\`ebre affine de l'espace
motivique des chemins de $a$ \`a $b$.
Pour trois points $a$, $b$, $c$, on dispose
aussi (loc. cit.)d'un coproduit
$$
{}_c A_a(X)(X)\to {}_cA_b(X)\otimes
{}_bA_a(X)
$$
correspondant \`a la composition des chemins.
Pour $a=b$, $\Spec({}_aA_a(X))$ est le {\it
groupe fondamental rendu unipotent motivique}.

\subhead
3.13
\endsubhead
Prenons pour $X$ le compl\'ement, dans une
droite projective $P$ sur $K$, d'un ensemble
fini de points rationnels, et montrons comment
comparer ces constructions \`a celles de
Deligne (1989).
En r\'ealisation de de Rham, $X$ fournit le
complexe (\`a diff\'erentielle nulle) des
formes diff\'erentielles \`a p\^{o}les
logarithmiques:
$$
K\mapright{0}\Omega_P^1(\log\,S).\tag3.13.1
$$
Pour $X^n$, on obtient la uissance tensorielle
$n^{\ieme}$ de (3.13.1).
Ces puissances tensorielles forment un
syst\`eme simplicial de complexes.
Normalisant, ce qui revient \`a supprimer $K$
dans (3.13.1) (cf. (3.11.1)), on obtient
apr\`es troncation en $n$
$$
\oplusop_{m\Le n}\Omega_p^1(\log\,S)^{\otimes m}.
\tag3.13.2
$$

La comparaison avec la r\'ealisation de Betti
est par les int\'egrales it\'er\'ees de Chen,
et ceci est compatible aux constructions de
Deligne (1989).
Quant \`a la r\'ealisation en cohomologie
$\ell$-adique, les arguments prouvant 3.4
s'appliquent aussi bien dans le cadre
$\ell$-adique.


